% =====================================================================
%  Harald Høffding, Herbert Spencer's Philosophy
%  Translated from the Danish: „Herbert Spencers Filosofi“, Det nittende Aarhundrede (April--September 1875), May 1875, pp. 106--130.
%
%  TRANSLATION -- GENERATED; DO NOT EDIT BY HAND.  Built by ebook-method/hoeffding-1875-76/tr_parts.py build spencer-1875 --write
%  from the parts in ebook-method/hoeffding-1875-76/tr/ (GLOSSARY.md, TRANSLATOR-BRIEF.md).  Made from the collated
%  transcription texts/hoeffding/spencer-1875/transcription.tex (2026-10-15); [N] in the margin = the page of the original.
%  Terms: GLOSSARY.md (choices marked ◆ there for Hans to review).
% =====================================================================
\documentclass[12pt,a4paper,openany]{book}
\usepackage[utf8]{inputenc}
\usepackage[T1]{fontenc}
\usepackage{libertinus}
\usepackage{libertinust1math}
\usepackage[german,french,english]{babel}
\usepackage[protrusion=true,expansion=true]{microtype}
\usepackage{setspace}
\usepackage[margin=1.3in]{geometry}
\usepackage{fancyhdr}
\usepackage{hyperref}
\hypersetup{pdftitle={Herbert Spencer's Philosophy}, pdfauthor={Harald Høffding}, hidelinks}
\pagestyle{fancy}
\fancyhf{}
\fancyhead[LE,RO]{\thepage}
\fancyhead[C]{\textit{H. Høffding: Herbert Spencer's Philosophy (1875)}}
\fancypagestyle{plain}{\fancyhf{}\fancyhead[LE,RO]{\thepage}}
\setstretch{1.3}
\usepackage{marginnote}
\usepackage{graphicx}
\DeclareUnicodeCharacter{0254}{\reflectbox{c}}   % ɔ (id est)
\renewcommand{\thefootnote}{\ifcase\value{footnote}\or *\or **\or ***\else\arabic{footnote}\fi)}
\newcommand{\opage}[1]{\marginnote{\footnotesize\textit{[#1]}}}
\newcommand{\signature}[1]{\par\smallskip\hfill #1\hspace*{2em}\par}
\pdfstringdefDisableCommands{\renewcommand{\footnote}[2][]{}\renewcommand{\opage}[1]{}}

\begin{document}

\begin{titlepage}
\begin{center}
{\LARGE\bfseries The Nineteenth Century.}\\[0.6em]{\small A Monthly for Literature and Criticism}\\[0.4em]{\small edited by}\\[0.2em]{\large Georg Brandes and Edvard Brandes.}\\[0.8em]{\small 1875. April--September.}\\\vfill{\small Copenhagen. Published by Gyldendal (F. Hegel). 1875.}\\[2em]
{\small Translated from the Danish}
\end{center}
\end{titlepage}
\clearpage

% --- p. 106 ---
\setcounter{footnote}{0}%
\opage{106}\begin{center}{\Large\sffamily\bfseries Herbert Spencer's Philosophy.\footnote[1]{In essentials a lecture, delivered at the Students' Association on 13 February 1875.}}\\[0.5em]\rule{4em}{0.4pt}\end{center}\medskip

When a scientific endeavor is called
philosophical, there may be two apparently
very different characteristics that entitle it
to this name. Either it is because it has a
decidedly critical character, turns doubtingly and probingly
against previously established presuppositions and results,
or because it seeks to gather all the results and points of view
of knowledge into one fundamental view, and thus
by virtue of its systematic character. Yet
these qualities do not exclude one another; to a certain
degree they are found united in every thinker, only that according to his
character and the age in which he lives,
either criticism or system will predominate,
--- and the ideal perfection of philosophy would be a thoroughgoing
criticism and an all-embracing system. It is also
easy to show that criticism and system point
to one another reciprocally. Criticism is distinction, analysis. It
cannot rest, then, until it has found all the elements,
all the sides in the nature of the objects with which it is occupied.
But neither can thought settle down to
rest in what is separated and incoherent; the results of criticism
and analysis are therefore used in turn as stones
% --- p. 107 ---
\setcounter{footnote}{0}%
\opage{107}for a new building, which is the firmer and surer the deeper
and more thorough the preceding criticism has been. These
two different tendencies are for thinking what inhalation
and exhalation are for the organism, both equally
indispensable, and on their rhythmic alternation the health of thought
depends.

In general it will surely be thought that in our
time the movement goes predominantly in the critical direction.
Our time, uncertain as it is of itself in theory as
in practice, a period of doubt and pessimism, seems not
favorable to systems and syntheses. And yet there may be,
in secret, through contributions from many different sides,
a store of stones gathered for a building which some day --- however
it may go with the provisional sketches --- will
be able to rise as a mighty temple, an expression of
all the ideal strivings of the human mind.

Whoever compares the nineteenth century with
its predecessor will not be able to deny that they stand
to one another as system to criticism. The history of the eighteenth century
is a process of dissolution; the nineteenth
strives, in many different forms, after organization.
As far as philosophy in particular is concerned, the most important attempts
in the eighteenth century --- Locke, Condillac, and Kant ---
have a decidedly critical and analytic character, whereas the
systematic tendency in our century is likewise represented
by the most eminent thinkers among the three
principal nations: Hegel, Comte, and Spencer. We find here
one and the same fundamental task treated in three different
ways.

For Hegel, who takes his stand on the standpoint of a priori speculation,
it is the life of thought that pulsates in
everything. The stages that thought passes through in order to reach its
perfection in the concept are one with those through which
all reality unfolds: the two movements coincide.
Although Hegel starts from the purely logical,
he has nevertheless seized the fruitful point of view of seeing this
% --- p. 108 ---
\setcounter{footnote}{0}%
\opage{108}as movement, process, development; thereby it became possible for him
to attempt a systematization of the whole content of
knowledge. But even if he was right in his conception
of the process of thought, he was nevertheless not right in
transferring without further ado the laws valid for this process
to the real.

Comte sought to free himself from all speculation and
took his stand on the standpoint of positive science.
He wished to explain man by the world, not the world
by man. And yet he is related to Hegel
by his decided systematic tendency, his striving for
and sense of totality.\footnote[1]{Hegel and Comte expressed mutual sympathy for each other's endeavors, admittedly without exhaustive knowledge of them. See Comte's correspondence with Eichthal in Littré: Auguste Comte et la philosophie positive.} He wished to abolish the isolation of the individual
sciences by gathering their principles and results in a general classification,
in order to show how
they finally merge into the one, all-embracing science:
the science of humanity. His fault becomes the
opposite of Hegel's: with the latter the manifoldness of reality was
violently pressed under the unity of thought; with
Comte the unity of thought is lost in the manifoldness of reality,
for mere classification cannot bring about
a real unity.

Spencer's attempt, with which we are here to occupy
ourselves more closely, stands in a way midway between Hegel's
and Comte's. Philosophy, he says, is knowledge developed into unity and
wholeness. We have thus not reached the
goal as long as the individual sciences and their content
still, as with Comte, form a series; we
must find one common law, one principle, in which all that we
know finds its expression. But this principle Spencer
does not find, like Hegel, immediately in the very nature of thought.
He finds it by the way of induction, by
% --- p. 109 ---
\setcounter{footnote}{0}%
\opage{109}comparison of phenomena in all the different domains,
just as Kepler in his day found his famous laws by
observing the planet Mars at different points of
its orbit. But what thus proves to be a
general law, valid for all phenomena, must then
also prove to stand in the most intimate connection with
the ultimate principles of our knowledge as a whole.

We understand a phenomenon only when we can follow
it from its first beginning and through the various stages
of its development. Just as knowledge of a people means
knowledge of its history, knowledge of a man,
knowledge of his family, childhood, youth, and manhood,
so we also know a tree or an animal only when we
know how they gradually unfold from the
invisible germs, and a chemical substance when we know of what
elements it is composed. Now it is clear that when
a phenomenon arises, this means that elements which
before existed in a scattered state are gathered and combined
in a peculiar way. Water arises through a certain
definite combination of oxygen and hydrogen atoms, the plant and
the animal through organic assimilation of certain chemical substances,
a man's character through the interaction between his
innate and inherited nature and the external conditions, etc.
Every phenomenon thus arises through a process of concentration.
The individual elements lose their independent
and autonomous motion and are brought for a time into equilibrium
in a certain definite way. But only for a
time. For as every phenomenon passes through a process of evolution,
so it also passes through a dissolution.
Concentration brings the parts into equilibrium; but this
does not last, it is superseded by a scattering of the elements
that for a time had combined to form
a peculiar whole. The history of every phenomenon
consists in a process of evolution and a process of dissolution.

Yet we can characterize these processes somewhat more closely. When evolution is a concentration of
% --- p. 110 ---
\setcounter{footnote}{0}%
\opage{110}previously scattered elements, then it is clear in the first place
that it denotes a transition from a less coherent
to a more coherent state. According to
the Kant-Laplacean hypothesis our solar system has thus
evolved by the original nebular mass
contracting to form individual globes. The evolution
of the earth consists in a condensation. Organic
life rests on the close coherence of the elements
that are taken up from the surrounding world through nutrition.
Human society is at its lowest stage a
horde of savages, which not only lacks inner coherence
but is also isolated from other societies; the further evolution
advances, the more coherence there is within
the single people and among all the single peoples mutually.

That which is thus brought into closer coherence in the course of evolution
passes over, in the second place, from a
homogeneous to a heterogeneous state. In the original nebular mass
there were no essential differences; these arise
only with the concentration of the individual globes. In the
evolution of the earth an opposition first forms
between crust and core, then on the crust itself an opposition
between mountains and plains, seas and continents.
In the evolution of the organism as well as of society
there sets in a division of labor, a thoroughgoing
specialization of tasks and functions.

In the third place, evolution consists in a transition
from the indefinite to the definite. Order everywhere supersedes
chaos; pronounced individual forms take shape:
thus the different globes, the different nations,
estates, and personalities.

These three laws, different forms of the general
law of evolution, are the main points of view to which every
consideration of phenomena leads us. They are the
general generalizations of all our experience. But
the question now arises whether a deeper foundation
and derivation of these laws might not be possible,
% --- p. 111 ---
\setcounter{footnote}{0}%
\opage{111}just as the Keplerian laws first received their proper foundation
when Newton traced them back to the law of gravitation.

Every phenomenon we necessarily refer to some
earlier phenomenon; in this consists, as we have seen,
the explanation we can attain of phenomena; explanation
of a phenomenon is the demonstration of its development.
The one link in the infinite series of phenomena illuminates the
other. The fundamental activity of consciousness is a constant
transition from one sensation or idea to
another. But then it is also clear that our consciousness
can function only, or in other words, that we can only
think, so long as such a transition from one link to
another is possible. A phenomenon of such a nature that
in its appearance and in its whole being it entirely fails
to cohere with or point to anything whatever else,
is an impossibility that we cannot even construct for
our imagination. The fundamental presupposition of thinking is therefore
that there is no absolute beginning or end, ---
no transition from nothing to something or from something to
nothing: such a transition would indeed be self-contradictory,
would be a relation with only one term!

If an absolute coming-to-be and passing-away is impossible, then
the sum of matter must also always be the same, and
motion must be continuous. Matter we know
only through the resistance it offers to our movement,
and motion only through the change it brings about
in our consciousness: the concept of force thus lies
at the basis of our concepts of matter and motion,
and the persistence of force is therefore the fundamental presupposition
on which all knowledge rests. The proposition that force
persists is the only one that goes beyond all experience,
because it lies at the basis of all experience.

Modern chemistry and physics have indeed proved that the
material elements never become nothing, but at every
disappearance of a body merely pass over into new combinations,
and that no force becomes nothing, but through
% --- p. 112 ---
\setcounter{footnote}{0}%
\opage{112}its apparent disappearance is only converted into another
force. But the proposition of the persistence of force itself cannot
be proved inductively, for the simple reason that
every induction presupposes it. Every chemical and
physical experiment rests on this proposition: as soon as anything
is to be weighed or measured, a unit of force must be laid
at the basis; but if this unit of force could change,
the result of the experiment would be illusory. All
scientific research therefore proceeds from this presupposition.

The deepest foundation, then, that we can obtain for
any empirical law consists in tracing it back to the proposition
of the persistence of force. Let us then look again at
the laws of evolution of phenomena. Originally the elements are
scattered in a homogeneous mass. But since the individual
parts of this receive different influences from their surroundings,
they will gradually themselves also become
different. For if different forces acting from outside
had the same influence on the homogeneous elements,
this would conflict with the persistence of force:
the difference of the forces would not manifest itself
in a difference in the effect. And when
differences thus form within the originally
homogeneous mass, then on the other hand
the same force will have different effects on the different
parts. These two propositions: that different forces
exert different influence on homogeneous parts, and that the same
force exerts different influence on different parts, ---
both of which spring from the fundamental principle of the persistence
of force, --- contain the ground of the process of evolution
as a transition from indefinite and incoherent homogeneity
to definite, pronounced, and coherent heterogeneity.

The process of evolution does not now last forever. At last
a state is reached in which there is equilibrium between
% --- p. 113 ---
\setcounter{footnote}{0}%
\opage{113}the phenomenon and its surroundings. The forces acting from outside
no longer work toward evolution; the individual
elements have now lost their independent motion;
no further advance can then be made in the direction of concentration.
But equilibrium is only the beginning of dissolution:
for the external forces do not cease to act,
and their influence will now consist in the elements
being scattered and passing over again into independent motion. This too
follows directly from the persistence of force. When dissolution
has then reached its completion, external forces will
again be able to act concentratingly --- and so on ad
infinitum.

The life of the universe in all its parts and forms is thus
an infinite succession of evolution and dissolution.
This is the most general law to which experience can
lead us, and it finds, as we have seen, its foundation
in the fundamental law of our own consciousness.

But Spencer maintains most emphatically that we
do not thereby reach beyond the phenomenal and the relative.
The Absolute, that which is in and for itself, we can neither
imagine nor think: the activity of our consciousness
is indeed a transition from one term to another --- but the
Absolute, that which is complete and closed in itself, admits
of no such transition, since it has nothing outside
itself. Nevertheless we are compelled to assume an absolute principle:
otherwise the relative and conditioned itself would become absolute.
There must exist something other than mere limits and
conditions. We have moreover seen that beneath all the
changing of phenomena, their origin and dissolution,
there yet remains something --- force in itself, which can neither begin
nor cease. The Absolute, that which exists in itself,
is the force that stirs in everything, --- whose nature we neither
know nor can know, but which we must assume
as the ultimate ground and presupposition of everything that we
can know.

% --- p. 114 ---
\setcounter{footnote}{0}%
\opage{114}As science thus recognizes that, however
far it may reach in knowledge of the coherence of phenomena,
it will never be able to reach the deep source from which they
well forth, while yet we are of necessity led to
assume its existence, a union with the
religious consciousness becomes possible. The religions too maintain
that the core of existence is a mystery to men,
--- but they at the same time entangle themselves in the contradiction of claiming
to be in possession of very exact information
about this mystery. The fetish worshipper has the Absolute
before him in the block of wood or the stone, which he beats when
it will not hear him; at higher stages it is the animal or
human shape in which it is clothed; the history of the world
is still for many the autobiography of the absolute mind.
Against all positive forms of religion Fichte is right
when he says that a God we could comprehend would not
be a God, but an idol. Through the history of the
religions one can also trace a progressing consciousness
of this truth. But to this very day
the so-called pious, who cling to the traditional
finite expressions for the Unconditioned and, in order to
maintain these, harden themselves against the results of science,
thereby making themselves guilty of the worst unbelief: the fear
that truth might be harmful, --- are these
so-called pious less religious than the investigator, who
knows that all such forms are only symbols?

Also on this decisive point Spencer takes
a position that lies midway between Hegel and Comte.
Hegel finds in the speculative concepts, in which the core of existence
is exhaustively known, the truth of positive
religion contained as well. Comte, on the other hand, regarded
the religious standpoint as completely overcome; the
only religion that could now be in question was, according to
his conviction, the Religion of Humanity, which consists in
the heartfelt, even mystical devotion to humanity
and its evolution. Spencer maintains against this that
% --- p. 115 ---
\setcounter{footnote}{0}%
\opage{115}the religious feeling will always turn toward an absolute
ground of all existence, only that it must give up every
attempt to teach us anything about this absolute ground.
The universe itself in its richness and in its infinite rhythmic
motion is the only revelation of the eternal,
unknowable force that stirs in everything.

Although I wish here only to present Spencer's line of thought
and not my own position toward it, I cannot
withhold the remark that Spencer will hardly be able
to keep the concept of evolution wholly apart from all application
to the Absolute. When it is a fundamental law
of everything that it passes through evolution and dissolution, then
the core of existence cannot be entirely foreign to
this movement. Without wishing to venture upon new
speculative adventures, we can nevertheless, and we seem bound to,
assume that the rhythmic movement has significance
for the absolute principle itself. It is just as justified to
designate this as the principle of evolution as to use
the expression force about it; the law of evolution has indeed, according to
Spencer, the same universality and necessity as the concept
of force. Only under this presupposition does the
religious feeling also seem able to see the universe as a real
revelation of the Unconditioned; and only under this
presupposition can resignation to the laws of the world
become a source of new enthusiasm and striving.

Spencer himself has, moreover, given expression to this enthusiasm
in the following statement, called forth by the opposition
in which the investigator always feels himself to stand toward
the adherents of the positive religions, however much he may recognize,
or perhaps precisely because he recognizes, the relatively justified element
in their faith. ``He [the advanced thinker]
may, like any other man, rightly regard
himself as one of the countless organs through which the
unknown cause works, and when the unknown cause
produces a certain conviction in him, he is thereby authorized
to express it in speech and action.
% --- p. 116 ---
\setcounter{footnote}{0}%
\opage{116}For --- to apply the poet's words in their highest
sense ---
\begin{verse}\small
\hspace*{6em}Nature, after all, cannot\\
be improved by any means that Nature\\
has not itself created; and above this art,\\
which, as you say, aids Nature,\\
there is an art that Nature itself has created.
\end{verse}
\begin{flushright}\small\textbf{(Shakespeare: The Winter's Tale.)}\end{flushright}

Therefore the wise man will not regard the conviction
that stirs in him as accidental. He will without fear
utter the highest truth he can see, because he knows
that whatever the consequence may be, he thereby plays the right
role in the world, --- because he knows that if he \emph{can} bring about
the transformation he strives for, well! --- if he \emph{cannot}, also well, though not \emph{so} well! (First principles
3d ed. p. 123).'' Here Spencer evidently suggests another
relation between the absolute principle and the process
of evolution than in the determinations given above.

We return to the laws of evolution. It is in
the organic and the social domain that it most strikingly
shows itself. The periods of evolution are
here shorter, the conditions more easily surveyed. It is
here that the idea of evolution first dawned on thought, notably
by attention being drawn to the development of the sciences and of
social conditions through the ages and to
the growth of the individual organism. In his first more considerable
work, ``Social Statics'' (1850), Spencer already declared
that the same law held for organic and social
evolution. It is above all modern embryology, founded
by Harvey, Wolff, and v. Baer, that gave him
the impulse to one of his most important ideas: evolution as
a transition from a homogeneous to a heterogeneous state.
Step by step he was then led to apply this
law to the universe as a whole.

Spencer was thus fully prepared to enter upon
the grand hypothesis that Darwin put forward in his
% --- p. 117 ---
\setcounter{footnote}{0}%
\opage{117}famous work on the origin of species (1859). Darwin
also names him here as one of his predecessors. But
Spencer in this question does not place the chief weight
on the same factor as Darwin. The latter regards
the distinct species as having arisen, in the struggle for life continued through countless
generations, in that in
every generation the individuals that are not suited and do not
have the strength to maintain their existence in the face of their surroundings
succumb, are displaced by the stronger or
more fit, and are excluded from reproduction. Since the race
thus continually propagates itself through the individuals most fit
to live under the given conditions,
a form will gradually evolve that is in as
good harmony with the conditions as possible. This is
what Darwin calls natural selection or the selection of quality.
According to Spencer Darwin has laid too much
weight on this indirect evolution. He himself takes his
starting point from the experience that exercise, function,
develops the organ. When a muscle is frequently used, it needs
more nourishment; more blood then flows to
it than before, and it will thereby increase in size. Now
there is indeed in the single individual a definite limit to
this increase; a point comes where the blood vessels
are not large enough to yield the required increased
supply. But it is also a fact that a structure developed by
increased function can be inherited by
the descendants. One generation then continues what the
preceding has begun, and thus the organization can
step by step be developed into such a form as the
external conditions require. Function may \emph{in this respect} be said
to precede structure. The lowest organisms, e.g.
the Rhizopods, are thus examples of life and functioning
without real organization\footnote[1]{It is of course not Spencer's meaning that function can be independent of structure. In that case G. H. \emph{Leves} \opage{118}would be right in his polemic against the proposition that function precedes structure. (Fortnigthtly Review. New Serie, vol. IV p. 69---73). What Spencer means is surely only that function can become the cause of the development of new structural relations, which agrees directly with experience. See his ``Principles of biology'', I p. 185---199.}.

% --- p. 118 ---
\setcounter{footnote}{0}%
\opage{118}Spencer admits that the indirect process of evolution,
which consists in the perishing of the imperfect individuals and
the propagation of the race through the more perfect, plays a
great part, especially the lower the organism stands. But the
more developed the organic form is, the more will the direct
process of evolution, by which one generation hands over to
the next the results of its struggles and experiences,
deposited in the organization, take the place of the indirect.
This happens in particular when a nervous and muscular system has
developed, and with it the capacity to vary actions
according to circumstances. When the organism is strong enough,
it will be able to adjust itself; only when it is too weak does
it perish.

According to Spencer the choice here lies only between the theory of evolution
and the theory of creation. Either the species have developed
according to natural laws under the interaction of natural
conditions, or else each species has been created separately after a
fixed type. To leave the whole question open cannot
be done in the long run; and when one in one's view of the world at all
relies on natural causes, the choice
between evolution and creation cannot be difficult. Creation
is even in the last instance unthinkable, cannot even be
constructed for the imagination: it expresses a transition from
nothing to something, --- which, as we have already shown,
conflicts with the fundamental law of our thought and consciousness.
One may therefore say that no one in reality believes in
a creation; for in order to believe one must at least
really be able to imagine what one is to assume, but
% --- p. 119 ---
\setcounter{footnote}{0}%
\opage{119}that is here impossible. One believes that one believes\footnote[1]{Cf.\ \emph{Feuerbach's} epigram ``Potenzirter Glaube'': Auf der höchsten Potenz steht jetzt in den Frommen der Glaube; Denn es wird ja sogar selber der Glaube geglaubt.}. And how
many particular difficulties the hypothesis of creation would
not lead to! The Creator would have had to hold fast to a type,
even where it was unsuitable; the seal, e.g., would have
been created for its present habitat and yet would
have the limbs of a mammal, which under the hypothesis of evolution
becomes quite explicable, when one thinks of the seal as
a land animal that has gradually accommodated itself to
life in the sea. It would, finally, follow directly from the hypothesis of creation
that certain species were destined to inflict
pain and destruction on others. The hypothesis of evolution, which does not at all
enter upon the transcendent cause, sees in the pain and struggle
of organic life one of the most important driving wheels
that lead forward to greater harmony; it certainly
cannot --- any more than any other doctrine
can --- tell us why pain and want are necessary
driving wheels in the life of the world; but it relieves us of the
belief that they should have been produced by design\footnote[2]{On the occasion of the sharp opposition Spencer maintains between evolution and creation, he has had a little polemic with Gladstone. The famous statesman and scholar was willing enough to go along with the general laws and the natural causes, but he reproached science, and especially the hypothesis of evolution, with ``entering into strife with Providence''; in particular he expressed the conviction that Providence in due time caused the great men to arise, and denied that this could be explained by natural laws alone. Spencer compares Gladstone's polemic against natural explanation with the Greeks' indignation at the physical conception of the heavenly bodies and their motions. Especially in the domain of social and personal life the natural and the theological conceptions run side by side in our time, without the contradiction being noticed. ``We proved our national skill in entertaining ideas that mutually cancel each other, when an unpopular prince suddenly won popularity by \opage{120}passing through certain morbid changes in his blood, and when on the occasion of his recovery we at one and the same time acknowledged the help of Providence and the working of natural causes by a thanksgiving to God and a baronetcy for the physician.'' \emph{Spencer}: The study of sociology. p. 392---394; 423---426.}.

% --- p. 120 ---
\setcounter{footnote}{0}%
\opage{120}The hypothesis of evolution brings, moreover, as we have already
seen, all the domains of our knowledge into a magnificent
coherence. The history of the solar system, of the globe, of the organism,
of language, of science and of human society falls under one and the same general
law. Under this general law the phenomena
of the life of consciousness also fall. Psychology is only a part of
biology, which in turn is a detached part of geology,
just as this again is a special part of astronomy.
Our distinctions and boundaries are subjective and arbitrary
in face of the infinite coherence of the universe. All is
in continuity. Thus one can therefore not point out
any definite point where the conscious life separates itself
from the unconscious. The transition from the physiological to the
psychological phenomena takes place through the increasing concentration
of the nervous system, whereby the scattered influences
are gathered and brought into connection with one another.
Hereby a higher stage of the general life-process
becomes possible: the impressions can now be compared and
appraised, and a more composite and more adequate reaction
thereby becomes possible on the part of the organism. The life of consciousness
develops step by step from the lowest forms
to the highest. In the lowest animals only dim sensation
and next to no combination and comparison
of the particular sensations is possible. It is, accordingly,
only very one-sided and limited relations of interaction
with the external world that can here take place.
Think, e.g., of the life of consciousness in an intestinal worm
or an oyster. The richer the organization becomes,
the more exact, many-sided and comprehensive becomes
% --- p. 121 ---
\setcounter{footnote}{0}%
\opage{121}the interaction. The human mind penetrates
the globe and the universe and has besides in itself a rich
and manifold world. Compare the consciousness of that intestinal worm
with Newton's, who propounds the laws of the universe,
or with Shakespeare's, who probes the soul's deepest
drives and thoughts. And yet the psychology of animals and of man
shows us a continuous series of intermediate stages
between these opposites.

It is the same within the human mind
itself. The various so-called faculties
--- instinct, reason, observation, the concepts, memory
etc. --- designate only different degrees or forms of
one and the same process. From the simplest combination
of sensations or presentations up to the most abstract
and intricate scientific reasoning it is
one and the same fundamental activity that manifests itself: namely,
observation of likeness and unlikeness between impressions or
ideas. All thinking is a classification, and
the development of mental life depends on how comprehensive, composite
and exact this classification is. The highest
form of thought, the quantitative or mathematical,
owes its advantages over the other forms to the circumstance
that it presupposes homogeneity in its
objects and investigates only the relation between their
magnitudes; but it consists in itself just as well in an
observation of likeness and unlikeness as the most elementary
sense-perception.

The prevailing tendency within English philosophy,
which in our day has had an outstanding representative
in John Stuart Mill, was decidedly empiricism, i.e.
it derived all knowledge from experience and conceived
consciousness as a blank slate (tabula rasa), which was gradually
written upon by impressions, which in turn were combined into
ideas and concepts. Against this conception Spencer
declares himself. If the mind, he says, prior to
all experience is a tabula rasa, where does the capacity
% --- p. 122 ---
\setcounter{footnote}{0}%
\opage{122}to organize and work up the experiences come from? and whence
come the differences we find in this respect
among different races and individuals? Why is a nervous system
needed at all? But the most important point
is that thoroughgoing empiricism would abolish the coherence and necessity
of knowledge. If it is to be possible at all
to learn from experience, we must
have a principle, a criterion, by which we can
measure it, and by virtue of which every progress in
knowledge is made. This last criterion Spencer finds
in the impossibility of the opposite. What we cannot
think, what contains a self-contradiction, we cannot
accept as true either. This is the test to which
we are referred back in every debate, and without which
empiricism itself could not prove a single one of its
propositions.

Although Spencer thus grants Leibniz and Kant that they are right
in their polemic against empiricism, he considers it just as
impossible to regard the forms of knowledge as fixed,
implanted without influence from the surrounding
world. The error of empiricism is that it will derive everything from
\emph{individual} experience. The insufficiency of this
points us beyond the individual to the race. The single individual
does not begin absolutely from the start; it has in its innate
organization the results of the experiences of all earlier generations,
in so far as they have passed into the
structure. The universal and pervading forms in
things must gradually produce corresponding forms
in the organisms that in infinite series of ages have been subjected
to their influence. The original forms of knowledge
or the accumulated and transmitted structural modifications
that are produced by experience, lie hidden in every
newborn individual, stir in an indeterminate way
at the first individual experiences and become gradually more determinately
marked as such individual experiences
are multiplied.

% --- p. 123 ---
\setcounter{footnote}{0}%
\opage{123}Spencer thus arrives at the result: a posteriori
for the race, but a priori for the individual; what is originally
experience becomes, through the organization,
principle. In this way he seeks to get beyond
the opposition between the two main philosophical tendencies in
modern times, beyond the strife between Locke and Kant.

Yet he still counts himself among Locke's school,
because according to the hypothesis of evolution all knowledge stems
from experience, only that it is the experience of the race, not of the individual.
To this I must remark that Spencer's criticism
of empiricism holds just as well when it is a question of
the experience of the race as when it is a question of the individual's.
In order that the race may have experiences, it must have an
organization. Here an essential relation is in question,
which is the same at every stage. Moreover both
Kant and Locke treated the question exclusively
from the standpoint of the individual. When Spencer here grants Kant
that he is right, he stands evidently, with regard to the problem as a
whole, nearer to the German philosopher than he stands to his
English predecessors. Kant, after all, also teaches expressly
that all our knowledge begins with experience; he denies only
that it wholly and entirely arises from experience\footnote[1]{In my more detailed exposition of Spencer's doctrine in the work ``Den engelske Filosofi i vor Tid'' I have already pointed out that Spencer, despite his own declaration, must be said to stand nearer to Kant than to Locke (p. 170 f.). From the later published 3rd volume of Spencer's ``Essays'' I see that \emph{Max Müller} has made a similar remark. (p. 275: ``In admitting that there is something in our mind, which is not the result of \emph{our own} à posteriori experience, Mr. Spencer is a thorough Kantian'').}.
Spencer introduces a wholly new standpoint in the treatment of the problem of knowledge;
therefore it is difficult or
impossible to refer him to any of the tendencies into
which earlier thought has split in the debate
on this problem. I cannot see otherwise than
% --- p. 124 ---
\setcounter{footnote}{0}%
\opage{124}that by this new standpoint he has brought the whole question
a considerable step forward\footnote[1]{From an interesting letter from \emph{Stuart Mill} of 29 January 1872 to the physiologist \emph{Carpenter}, which is quoted in the latter's ``Mental Fysiology'' p. 486, one sees that the greatest representative of pure empiricism in our time, toward the end of his life, approached a view similar to Spencer's.}.

If it thus holds of man's knowledge that
it becomes intelligible only when one goes beyond the
single individual and sees him in his intimate connection
with the whole race, it nevertheless finds application in still higher
degree when it is a question of drive, feeling and
will. Drive and feeling stir in the individual before he
has made any experience that could produce them;
this holds in particular of the strongest of all drives,
the sexual drive. The ultimate motives of human
action must be sought in the fundamental drives that the race has
acquired in the course of its development. No insight, however clear,
compels action except by calling forth a
sufficiently strong and lasting feeling. It is thus not
true that ideas rule the world. Ideas serve only as
leaders; it is the individual and national character
that determines whether the leader is followed. Ideas come
into being mainly through the work of individuals and therefore cannot
directly overcome the deep and mighty impulses
that propagate themselves down through the race. It follows also
of itself that no ideas can arise that
stand in complete opposition to the prevailing character.
When higher ideas have developed, they certainly
contribute to further development and advance. But the most important
factor is and remains the very adjustment to the
natural and social conditions that every generation must
attempt. Character is developed only by action, and
action is called forth by the real conditions. True
education --- within the family as within the nation
% --- p. 125 ---
\setcounter{footnote}{0}%
\opage{125}as a whole --- must therefore consist in this, that individuals are kept
in interaction with the actual conditions.

But this insight is, according to Spencer, very far from
being general. Of all the prejudices that hinder
sound thinking in the domain of education and social life,
none is more widespread and more harmful than the one
that parents and governments can accomplish something positive.
By interfering directly they only hinder the natural
interaction between individuals and their surroundings.
They have only to see to it that the necessary conditions
are present for the individual's existence and development. The
presupposition from which one generally proceeds is
that the virtues are found in the rulers, the vices in the
ruled. One looks solely to the \emph{ordering} force
and forgets that what is essential must nevertheless always be
the very thing that is \emph{ordered}.

The modern reformers, who think they can save
humanity by their systems of instruction or
constitutions of the state, are entangled in no less error
than the conservatives and reactionaries; both parties have
this in common, that they would bring individuals under tutelage.
It is a sort of fetish worship of the governing power
that has developed, and with which radicals and
Tories are equally infected. Men cannot, after all,
produce new forces; they must work with the
means that Nature gives. The character of society rests on
the character of the individuals by whose union it exists.
But this character is not changed all at once
by being trained in certain intellectual skills\footnote[1]{``Ignorance and crime are not cause and effect, but effects of one and the same cause. To be \emph{wholly} without knowledge means to have lived among people who have the strongest impulses to crime; to have received \emph{some} instruction means to belong to a class that is subject to less strong temptations; to be \emph{well} instructed means to have lived under conditions where the usual motives to \opage{126}crime hardly stir.'' \emph{Spencer}: Social Statics. p. 379 f.} or
% --- p. 126 ---
\setcounter{footnote}{0}%
\opage{126}by being brought into a certain system fetched from outside. Just
as little as one can escape the influence of gravity
by placing one's limbs in a certain definite position,
just as little can men hope to escape the
wickedness that weighs them down by arranging themselves according to this
or that political pattern. One constantly overlooks the
hidden forces working in the depths and stares at the forms
and factors that move on the surface.

Primitive man had to have wild and egoistic instincts
in order to be able to survive in the hard struggle against Nature,
the animals and his equals. When a society began to form,
--- an important means for carrying on this same struggle,
--- certain bonds were imposed on individuals that were
necessary if society as a whole was to exist.
That society and those individuals advanced furthest and survived best
in the struggle among whom the adjustment to the demands of
social life proceeded furthest. The egoistic instincts
therefore did not long remain sole rulers; the sympathetic
drives, which even from the first have so deep a ground
in the relation of the sexes and in the relation between parents and children,
had to increase more and more in strength and extent. The
one respects the right of the other and perhaps even works
unselfishly for his welfare. Egoism thus passes over
into altruism, to use Comte's apt expression, and
the ideal of social development, which also contains the
ethical ideal, is a condition in which all work together,
each according to his individuality and with as much freedom
as can exist with the freedom of all others.

But this is a condition to which only natural
development can lead. The government is to see only to
security and justice, the fundamental conditions for
development to be able to proceed. Certainly in society
as in every organism a certain organization, a certain
% --- p. 127 ---
\setcounter{footnote}{0}%
\opage{127}structure is necessary. But the less this proceeds from
the very nature and interaction of the elements, the more
artificial and imposed from above it is, the greater the obstacle that
progressive development will have to overcome
through it. The development of language and economic development,
the producing, exchanging and distributing processes,
are examples of results that have been reached by a natural
route, where each individual, working under the given
conditions and according to natural ends, produced
results that lay far beyond the individual horizon
both in extent and in significance. As
the egoistic instincts are pushed back, unselfish
endeavors on the part of individuals, or free associations
of individuals, will also throw their weight into the scale and accomplish
what the former could not. The state, which cannot even
manage its principal task, an impartial and energetic
upholding of right and security, can still less accomplish
anything in other domains and will in any case do more
harm than good. Humanity is in a state of transition;
it is to lay aside its original robber nature and
become fitted for perfect social life. And this
state of transition must necessarily be unhappy: all
unhappiness consists in disagreement between the inner
organization and the external conditions. But the sufferings
must be gone through. They can be alleviated by sympathy, by
the love that helps others to help themselves;
but a certain measure of suffering is unavoidable. To
intervene violently would precisely increase and prolong the suffering.
There are larvae that live in the water and breathe by gills,
but in which lungs are already developing here for use
in a later period of the animal's life. Before the lungs
are fully developed, the animal cannot do without the gills;
one would kill it, if one suddenly wished to put
an end to the inconsistency of having lungs and gills
at once. And so it goes also in the
development of social life; we can retard it by directly
% --- p. 128 ---
\setcounter{footnote}{0}%
\opage{128}
intervening, but we cannot hasten it in any other
way than by asserting the influence of the natural conditions.
And even in the sufferings that thereby necessarily
set in, he who fixes his gaze on development
and its goal as a whole will recognize a fighting love.

Spencer's ethical and sociological views are
not yet available in worked-out form. Of his philosophical system,
which is to comprise 10 volumes, exactly half
remains, which is to contain the principles of sociology and ethics\footnote[1]{The first five volumes are: First Principles, Principles of Biology (2 vols.) and Principles af Psychology (2 vols.).}.
But the general character and direction are clearly apparent
from what we know. The first point I will emphasize
is that Spencer aims at an altogether objective,
strictly causal conception of human social life.
He is not, as Stuart Mill is in all his social and
political essays, at the same time a reformer. This must
be kept well in view. He has been accused of fatalism,
because he lays so much weight on the fact that
human development is a natural development, and recently
a disciple of Stuart Mill has cited against him the following
words of the latter: ``Political institutions are the
work of men, owe their origin and
their whole existence to human will. Men do not wake up
one fine morning and find them grown up\footnote[2]{\emph{Cairnes}: Fortnightly Review. January 1875.}.'' But the question
is, according to Spencer, precisely whether men do not
in their very willing and acting obey a law of nature of which they
are not conscious. He wishes here again to cite Shakespeare's
words, that the art which is said to help Nature
has its ground in an art that Nature herself has produced.
And one will be forced to admit this, unless
one will assert a concept of freedom of the will that
bursts all lawfulness and coherence in the
psychological and social domain.

% --- p. 129 ---
\setcounter{footnote}{0}%
\opage{129}Spencer's doctrine will not push back conscious
striving and action either; it is only the way and direction
of this striving that the disagreement concerns. According to
his conviction the true and full development of
man's personal powers can take place only if one
does not prevent the individual from coming face to face
with the real conditions and taking up the struggle for existence
under the real and actual circumstances. The task
of the power of the state consists then for him, as for Wilhelm
v. Humboldt, in an emancipation of individual
freedom. On account of the conditions of development that are
required for the social life-process as for every life-process,
it is according to his conviction an illusion to
expect considerable results in short periods\footnote[1]{The famous embryologist \emph{Ernst} v. \emph{Baer}, from whom Spencer has learned so much, has come to a similar result by comparing social development with the organic. ``Es prägte sich mir tief die Ueberzeugung ein, dass die socialen Entwickelungen nicht anders vor sich gehen können als die organischen, nämlich allmählich, und nur dann gesunde Gestalt gewinnen, wenn die Bedingungen dazu vorher vollständig gegeben waren und längere Zeit bestanden hätten \dots{} Die Männer, welcke vorwärts wollen, --- und rückwärts zu wollen ist überhaupt gegen die Natur, --- würden gut thun, ihre Gedanken recht viel an die Beachtung der organischen Entwickelungen zu gewöhnen.'' Nachricht über Leben und Schriften K. E. v. Baer’s, mitgetheilt von ihm selbst. St. Petersburg 1865. p. 52. --- In his latest work ``The study of sociology'' Spencer demands as a necessary preparation for occupation with the science of society the study of the exact sciences, of biology and psychology.}.

But it cannot be denied that Spencer overlooks the
positive significance that the intervention of the power of the state may have
for the attainment of the end toward which social
development strives. Certainly he would hardly grant Buckle
that ``it would be to mock all sound reason,
if one ascribed to legislation any share
% --- p. 130 ---
\setcounter{footnote}{0}%
\opage{130}in progress or expected from future legislators
any other benefaction than the repeal of their
predecessors' work;'' --- certainly he declares (in the preface to the
second edition of ``Social Statics'') that he would now more than
before stress the relative value of political institutions;
--- but one nevertheless misses in him the clear
recognition of the educative significance of society. ``There
is,'' as I have expressed this in another place\footnote[1]{Den engelske Filosofi i vor Tid, p. 65.}, ``still
a long way to go before the truly natural development
can set in. The state, by whose help man
was led out of the animal state of nature, cannot let go of
him before it has led him to the human state of nature.''
It is on this that the relative justification of the socialist ideas
is grounded. Here too the solution of the
oppositions lies in the idea of evolution; the opposition
between individualism and socialism, which according to
Stuart Mill's statement\footnote[2]{John Stuart Mills Levnet fortalt af ham selv. Oversat af V. Arntzen, p. 259.} is diametrical, can be overcome only
by both tendencies being conceived as mutually conditioning
each other. But from such a solution we are still far
away. Spencer's sociology, which moves predominantly
in the theoretical domain, leads to a salutary skepticism
toward all provisional attempts at solution, while
it energetically points toward the most deeply working forces and
the one essential goal.
\signature{\textbf{Harald Høffding.}}

\end{document}
